\section{Формальная постановка задачи}

Математическое ядро работы состоит в задаче сравнения информационных режимов в линейной модели
в пространстве состояний, построенной на основе HANK/SSJ-откликов. HANK/SSJ-блок задаёт локальные
равновесные отклики экономики; разные наборы наблюдений превращаются в разные информационные
состояния регулятора; для каждого режима заново настраивается правило процентной ставки; затем
сравниваются оптимизированные потери. Поэтому основным объектом анализа является не отдельная
распределительная переменная, а переход между информационными режимами.

\subsection{Линеаризованная модель в пространстве состояний}

После перехода от HANK/SSJ-откликов к низкоразмерному состоянию в полной локальной постановке с
управлением динамику можно записать как
\[
q_{t+1}
=
Aq_t+B i_t+\varepsilon_{t+1},
\qquad
\mathbb E\varepsilon_{t+1}=0,
\qquad
\mathbb E\varepsilon_{t+1}\varepsilon_{t+1}^{\top}=\Sigma_\varepsilon .
\]
Здесь \(q_t\) включает агрегатные и распределительные компоненты,
\[
q_t=(\pi_t,y_t,C_t,\overline{\kappa}_t,\ell_t,\xi_t^r),
\]
\(i_t\) -- отклонение инструмента денежно-кредитной политики от стационарного значения, а
\(\varepsilon_{t+1}\) -- вектор экзогенных шоков. Матрицы перехода описывают поведение выбранного
вектора состояния в окрестности стационарного равновесия и используются для сравнения
информационных режимов.

В основной оценке на фиксированной траектории состояния правила сравниваются на общих тестовых
траекториях, порождённых HANK/SSJ-динамикой. Это позволяет сравнивать информационные состояния на одной и той же
экономической базе. Член \(B i_t\) явно используется в линейно-квадратичном ориентире и в проверке
с обратной связью ставки, где инструмент политики возвращается в динамику состояния.

\subsection{Информационное состояние регулятора}

Для каждого информационного режима \(\varphi\) регулятор получает сигнал
\[
o_t^\varphi
=
M_\varphi q_t+\nu_t^\varphi,
\qquad
\mathbb E\nu_t^\varphi=0,
\qquad
\mathbb E\nu_t^\varphi(\nu_t^\varphi)^{\top}=R_\varphi .
\]
Матрица \(M_\varphi\) задаёт, какие компоненты состояния доступны наблюдению, а \(R_\varphi\) описывает
точность соответствующих сигналов. Информационное множество регулятора:
\[
\mathcal I_t^\varphi
=
\sigma(o_0^\varphi,\ldots,o_t^\varphi,i_0,\ldots,i_{t-1}).
\]
Информационное состояние определяется как сжатие этой истории:
\[
s_t^\varphi
=
\Phi_\varphi(\mathcal I_t^\varphi).
\]
В разных режимах \(\Phi_\varphi\) может быть текущим наблюдением, набором лагов или фильтрованной
оценкой состояния. Тем самым распределительные статистики входят не как произвольные регрессоры,
а как элементы информационной структуры \(M_\varphi\) и соответствующего состояния \(s_t^\varphi\).

\subsection{Правило ставки и функция потерь}

Для каждого информационного режима \(\varphi\) задаётся класс правил
\[
\mathcal G_\varphi
=
\{g_\theta^\varphi:\ i_t=g_\theta^\varphi(s_t^\varphi),\ \theta\in\Theta_\varphi\}.
\]
В основной спецификации используется линейное правило с инерцией:
\[
i_t
=
\rho_i i_{t-1}+\theta^\top s_t^\varphi .
\]
Коэффициенты выбираются заново для каждого информационного режима:
\[
\theta_\varphi^\star
=
\arg\min_{\theta\in\Theta_\varphi}J(\theta,\varphi).
\]
Коэффициенты не переносятся между режимами: каждое информационное состояние оценивается при
собственной настройке правила в одном и том же классе.

Критерий является квадратичной функцией потерь стабилизации:
\[
L_t
=
\lambda_\pi \pi_t^2
+\lambda_y y_t^2
+\lambda_c C_t^2
+\lambda_i(i_t-i_{t-1})^2.
\]
Все макроэкономические переменные в функции потерь выражены как отклонения от стационарного
значения или целевого уровня. В частности, \(C_t\) понимается как отклонение совокупного
потребления от стационарного значения. Поэтому критерий стабилизирует инфляцию, выпуск,
потребление и сглаживание ставки. Конкретные значения весов функции потерь задаются в
расчётной спецификации и остаются неизменными во всех информационных режимах. Квадратичный
критерий используется как стабилизационная функция потерь: он фиксирует одинаковую шкалу
сравнения правил процентной ставки по инфляции, выпуску, потреблению и сглаживанию инструмента.

\subsection{Оптимизированные потери и предельная ценность информации}

Оптимизированные потери режима \(\varphi\) равны
\[
J_\varphi^\star
=
\min_{\theta\in\Theta_\varphi}
\mathbb E
\sum_{t=0}^{T}
\beta^t
L(q_t,i_t,i_{t-1}),
\qquad
i_t=g_\theta^\varphi(s_t^\varphi).
\]
Важна именно оптимизация внутри каждого информационного режима: сравниваются не фиксированные
коэффициенты правила и не прогнозные ошибки отдельных переменных, а минимальные потери внутри
одного класса правил при разных информационных режимах.

Предельная ценность перехода от информационного режима \(a\) к режиму \(b\) определяется как
\[
\Delta J(a \to b)
=
J_a^\star-J_b^\star .
\]
Положительное значение означает, что режим \(b\) позволяет снизить потери после того, как
правило процентной ставки заново оптимизировано под доступную информацию. Для главного сравнения работы
\[
\Delta J_{\mathrm{dist}}
=
J^\star_{\mathrm{agg}}
-
J^\star_{\mathrm{agg+dist}} .
\]
В этой постановке вопрос о пользе распределительных данных превращается в вопрос о снижении
оптимизированных потерь при расширении \(\mathcal I_t^\varphi\) и соответствующего состояния
\(s_t^\varphi\).

Для парных таблиц используется то же соглашение о знаке:
\[
\Delta J
=
J_{\mathrm{baseline}}-J_{\mathrm{alternative}}.
\]
Положительное значение \(\Delta J\) означает снижение потерь относительно базового информационного
режима.

\subsection{Фильтрация состояния}

В линейно-гауссовой постановке фильтрованная оценка состояния строится рекурсивно
\cite{kalman1960filtering}. Сначала
формируется прогноз:
\[
\widehat{q}_{t|t-1}
=
A\widehat{q}_{t-1|t-1}+B i_{t-1},
\]
\[
P_{t|t-1}
=
A P_{t-1|t-1}A^\top+\Sigma_\varepsilon .
\]
После наблюдения \(o_t^\varphi\) вычисляется калмановский коэффициент:
\[
K_t^\varphi
=
P_{t|t-1}M_\varphi^\top
\left(
M_\varphi P_{t|t-1}M_\varphi^\top+R_\varphi
\right)^{-1}.
\]
Оценка состояния обновляется как
\[
\widehat{q}_{t|t}
=
\widehat{q}_{t|t-1}
+
K_t^\varphi
\left(
o_t^\varphi-M_\varphi\widehat{q}_{t|t-1}
\right),
\]
а ковариация ошибки:
\[
P_{t|t}
=
(I-K_t^\varphi M_\varphi)P_{t|t-1}.
\]
Разные информационные состояния отличаются матрицей наблюдения \(M_\varphi\) и ковариацией ошибки
наблюдения \(R_\varphi\). Поэтому фильтрация позволяет сравнивать не только разные наборы переменных,
но и разные информационные структуры регулятора.

\subsection{Линейно-квадратичный ориентир}

Для сопоставления простых правил процентной ставки с решением той же линейной задачи используется
линейно-квадратичный ориентир. В режиме
полной информации решается LQR-задача для той же локальной системы. Поскольку функция потерь
содержит сглаживание ставки, удобно перейти к управлению изменением ставки:
\[
u_t=i_t-i_{t-1},
\qquad
\widetilde q_t=(q_t,i_{t-1}).
\]
Тогда расширенная система записывается как
\[
\widetilde q_{t+1}
=
\widetilde A\widetilde q_t+\widetilde B u_t+\widetilde\varepsilon_{t+1},
\]
а квадратичная функция потерь принимает стандартный вид
\[
L_t
=
\widetilde q_t^\top Q_L\widetilde q_t+u_t^\top R_Lu_t.
\]
Технически LQR-задача решается через обратную рекурсию Риккати для матрицы функции стоимости:
\[
\begin{aligned}
P_t
=\;&
Q_L+\beta \widetilde A^\top P_{t+1}\widetilde A\\
&-
\beta \widetilde A^\top P_{t+1}\widetilde B
\left(
R_L+\beta \widetilde B^\top P_{t+1}\widetilde B
\right)^{-1}
\beta \widetilde B^\top P_{t+1}\widetilde A .
\end{aligned}
\]
Соответствующее линейное правило для изменения ставки:
\[
u_t
=
-
\left(
R_L+\beta \widetilde B^\top P_{t+1}\widetilde B
\right)^{-1}
\beta \widetilde B^\top P_{t+1}\widetilde A\widetilde q_t .
\]
В режиме неполной информации используется принцип разделения: управляющая часть задаётся
LQR-рекурсией, а ненаблюдаемое состояние заменяется калмановской оценкой
\(\widehat{\widetilde{q}}_{t|t}\). Поэтому линейно-квадратичный ориентир показывает, какую
ценность распределительные наблюдения имеют в LQG-версии той же задачи.
